\documentclass[11pt]{article}
\usepackage[margin=1in]{geometry}
\usepackage{amsmath,amssymb,amsthm,longtable,booktabs,hyperref}
\newtheorem{theorem}{Theorem}[section]
\newtheorem{lemma}[theorem]{Lemma}
\newtheorem{proposition}[theorem]{Proposition}
\newtheorem{corollary}[theorem]{Corollary}
\emergencystretch=3em
\newcommand{\id}{\mathbb I}
\newcommand{\Herm}{\operatorname{Herm}}
\newcommand{\Sym}{\operatorname{Sym}}
\title{No four-fold unbiased measurements with four outcomes of rank two}
\author{Daniel Rupawalla\\ Independent Researcher\\ \texttt{danielrupawalla@utexas.edu}}
\date{10 October 2026}
\begin{document}
\maketitle
\begin{abstract}
We study the relation between spectral and algebraic higher-order
unbiasedness for equal-rank projective measurements. Exact averaged
fourth- and sixth-moment identities express deviations from the universal
spectral target as sums of squared operator defects. Consequently every
spectral $k$-fold unbiased tuple is pairwise mutually unbiased, and every
three-setting subtuple is an algebraic 3-fold unbiased triple (3-UM), at
arbitrary finite rank and outcome number. In particular there is no
three-outcome spectral $k$-UM for any $k\geq4$, extending the known
case $k=3$.
For four outcomes at rank two, a determinant-phase argument proves that
every 3-UM has a simultaneous quaternionic realization. Combined with the
moment identities, this shows that in this regime a spectral 4-UM is
automatically algebraic. We classify the possible third outcomes over every
pair that is not a direct sum of mutually unbiased bases, and show that
every 3-UM contains a pair that is such a direct sum. A uniform obstruction
for triples containing such a pair then shows that no four-fold unbiased
measurement, algebraic or spectral, with four outcomes of rank two exists.
Several steps are certified by exact symbolic computation, with scripts
provided; no numerical search is used in the proofs.
\end{abstract}

\section{Introduction and definitions}
Higher-order unbiasedness constrains quantum measurements beyond their
pairwise overlaps. Designolle and Farkas introduced algebraic and spectral
$k$-fold unbiased measurements and established their equivalence for pairs,
triples, rank-one measurements, and binary measurements \cite[Thm.~3]{DF}.
For four or more settings and more than two outcomes, the relation
between the two definitions is left open there, and no genuinely
higher-order example with $k>3$ and $n>2$ is known \cite[Sec.~VI]{DF}.
Their smallest nonbinary triple has four outcomes of rank two on
$\mathbb C^8$. They extend it to four settings whose every three-setting
subtuple is a 3-UM, and show that this quartet is neither an algebraic nor
a spectral 4-UM: its selected sums have fourth moments $64$ and $65$
rather than the target $257/4$ \cite[App.~C]{DF}. Whether any
four-outcome rank-two quartet can be four-fold unbiased is therefore the
natural existence question in this regime.

Our main results address these two questions in a precise order.
Theorem~\ref{thm:spectralheredity} proves that the spectral definition
forces pairwise and three-setting algebraic relations at every rank.
Designolle and Farkas observe that their four-outcome triple is the
complexification of three quaternionic mutually unbiased bases \cite{DF};
Theorem~\ref{thm:quaternionrigidity} shows that every four-outcome
rank-two 3-UM has this form.
Corollary~\ref{cor:unconditionalfour} uses this to show that a spectral
4-UM in this regime is automatically algebraic, which transfers the
algebraic obstruction below to the spectral definition.
Theorem~\ref{thm:commutingnonextension} excludes every extension of a
triple containing a commuting canonical pair. Theorem~\ref{thm:globalcommuting} proves that every strong triple
contains such a pair, and Corollary~\ref{cor:globalnonexistence}
therefore settles the four-setting existence question at rank two.
This is consistent with the quartet of \cite[App.~C]{DF}: its first two
settings are the computational and real-Hadamard measurements, which form
a commuting canonical pair.

The rank-one higher-order framework is due to Farkas \cite{F}.
Unitary-block representations and the canonical-form commutativity
criterion for higher-rank mutually unbiased measurements are developed in
\cite{FKN}; order-four
quaternionic and complex Hadamard classifications appear in \cite{CD,TZ}.
These supply ingredients for the structural and extension arguments.
Our arbitrary-rank statement concerns low-order spectral heredity, rather
than full algebraic--spectral equivalence at every rank and order.
The global nonexistence result concerns rank two; higher ranks require
a different structural argument.

\paragraph{Definitions.}
An $n$-outcome rank-$r$ projective measurement on $\mathbb C^{nr}$ is
an orthogonal resolution of the identity $(P_{a|x})_{a=1}^n$, with each
projector of rank $r$. For $k$ settings and an outcome string $\vec a$,
write $S_{\vec a}=\sum_{x=1}^kP_{a_x|x}$.
The tuple is a spectral $k$-UM when every selected sum has
characteristic polynomial $\mu_{k,n}(t)^r$, where
\[
 \mu_{k,n}(t)=\left(1-\frac1n\frac{d}{dt}\right)^kt^n.
\]
It is an algebraic $k$-UM when, for every selected outcome $A$ used
as anchor, the sum over all orderings of the other $k-1$ selected
outcomes satisfies
\[
 \sum_{\pi\in S_{k-1}} A B_{\pi(1)}\cdots B_{\pi(k-1)} A
 =\frac{(k-1)!}{n^{k-1}}A.
\]
We state both definitions for projective measurements of equal rank;
\cite{DF} define the algebraic notion for general measurements and show
that projectivity and equal rank follow \cite[Thm.~3(1)]{DF}.
The algebraic definition is hereditary \cite[Thm.~3(1)]{DF}. In particular, any algebraic
$k$-UM is pairwise mutually unbiased:
$PQP=P/n$ and $QPQ=Q/n$ for selected outcomes from distinct settings.
A \emph{strong triple} means an algebraic 3-UM; by \cite[Thm.~3(5)]{DF}
this is the same as a spectral 3-UM. At four outcomes its
anchored relation is
\begin{equation}\label{eq:triple}
 A(BC+CB)A=\frac18 A.
\end{equation}
For later use, $\Sym(B,C,X)$ denotes the sum of the six ordered products
of $B,C,X$.

For a rank-two mutually unbiased pair, choose computational coordinates
for the first measurement and write the second outcome frames as block
columns $(U_{ab}/2)_a$, with $U_{ab}\in U(2)$. Row and column frame
gauges set $U_{a0}=U_{0b}=\id_2$. We call the resulting blocks canonical,
and call the pair commuting when these blocks commute pairwise.
Our dephased blocks are a canonical form in the sense of \cite{FKN}.
By \cite[Prop.~2.8]{FKN}, in canonical form the blocks commute if and only
if the pair is a direct sum of mutually unbiased bases
\cite[Def.~2.7]{FKN}. Since that is a property of the pair itself, any two
dephased gauges of the same pair agree on whether the blocks commute.
Commutativity of the blocks does not mean that the outcome projectors of
the two measurements commute.

\section{Spectral rigidity from averaged moments}
\label{sec:spectralrigidity}
For $k$ rank-$r$, $n$-outcome PVMs on $\mathbb C^{nr}$, put
$S_{\vec a}=\sum_x P_{a_x|x}$, $h=1/n$, and let $\mathbb E$ denote uniform
averaging over the relevant outcome strings. Write $p_m(k,n)$ for the
sum of the $m$th powers of the zeros of $\mu_{k,n}$, including multiplicities.
For two settings define the nonnegative number
\[
 D_{xy}=\mathbb E_{a,b}\left\|P_{a|x}P_{b|y}P_{a|x}
                         -hP_{a|x}\right\|_{\mathrm{HS}}^2.
\]
For three distinct settings, write $P,Q,R$ for their selected outcomes and
define the three anchored defects
\[
 D_{xyz}^{P}=\mathbb E\left\|P(QR+RQ)P-2h^2P\right\|_{\mathrm{HS}}^2,
\]
with $D_{xyz}^{Q},D_{xyz}^{R}$ defined by exchanging the anchor.

\begin{lemma}[Exact averaged defect identities]
\label{lem:spectraldefects}
For arbitrary equal-rank PVMs,
\[
 \mathbb E\operatorname{Tr}S^4-rp_4(k,n)=2\sum_{x<y}D_{xy}.
 \tag{M4}
\]
If every pair is mutually unbiased, then
\[
 \mathbb E\operatorname{Tr}S^6-rp_6(k,n)
 =\sum_{x<y<z}(D_{xyz}^{P}+D_{xyz}^{Q}+D_{xyz}^{R}).
 \tag{M6}
\]
\end{lemma}
\begin{proof}
Expand the powers as words in selected projectors. Under the real trace,
cyclic rotations agree. Adjacent repeated letters collapse by projectivity.
A setting occurring just once can be averaged out, contributing $h\id$.
For (M6), pairwise mutual unbiasedness additionally replaces $PQP$ by $hP$.
The following table gives the constant contributions, divided by the
ambient dimension $nr$, for each specified subset of $\ell$ settings.
Multiply each row by $\binom{k}{\ell}$. The cycle column gives the
number of words left after these reductions.
\[
\begin{array}{c|cc|cc}
 &\multicolumn{2}{c|}{m=4}&\multicolumn{2}{c}{m=6}\rule{0pt}{13pt}\\
\ell&\text{constant}&\text{cycles}&\text{constant}&\text{cycles}\\\hline
1&h&0&h&0\\
2&12h^2&2&30h^2+30h^3+2h^4&0\\
3&36h^3&0&300h^3+234h^4&6\\
4&24h^4&0&1200h^4+360h^5&0\\
5&0&0&1800h^5&0\\
6&0&0&720h^6&0
\end{array}
\]
The two fourth-order cycle words are $PQPQ,QPQP$. Their average real
trace is $y=\mathbb E\operatorname{Tr}(PQPQ)$, and expansion of the
pair defect gives $D_{xy}=y-rh^2$. Replacing each cycle trace by $rh^2$
in the fourth-order expansion gives $rp_4(k,n)$.

The six sixth-order cycle words are $(PQR)^2$ and its cyclic rotations
and reversals. Put $x=\mathbb E\operatorname{Re}\operatorname{Tr}(PQR)^2$.
For $T=P Q R P$, pairwise compression gives
$\operatorname{Tr}TT^*=rh^3$, and completeness gives
$\mathbb E\operatorname{Re}\operatorname{Tr}T=rh^2$. Therefore
\[
 D_{xyz}^{P}=2x+2rh^3-4rh^4.
\]
Real-trace cyclicity gives exactly the same expression for the other
two anchors. Replacing each sixth-order cycle trace by
$r(2h^4-h^3)$ gives $rp_6(k,n)$, and its remaining excess is
$6[x-r(2h^4-h^3)]=D_{xyz}^{P}+D_{xyz}^{Q}+D_{xyz}^{R}$.

For completeness, the Newton identities needed for both replacements are
\begin{align*}
 h p_4(k,n)&=kh+\binom{k}{2}(12h^2+2h^3)
                  +36\binom{k}{3}h^3+24\binom{k}{4}h^4,\\
 h p_6(k,n)&=kh+\binom{k}{2}(30h^2+30h^3+2h^4)\\
 &\quad+\binom{k}{3}(300h^3+228h^4+12h^5)
 +\binom{k}{4}(1200h^4+360h^5)\\
 &\quad+1800\binom{k}{5}h^5+720\binom{k}{6}h^6.
\end{align*}
These follow by substituting the coefficients of $\mu_{k,n}$ into
Newton's recurrence. They hold also when $n<6$ or $k<6$, since the
corresponding falling factorials vanish. The accompanying verifier
checks the recurrences as polynomial identities in $k,h$, rather than
interpolating sampled parameter values, and enumerates all word types
in the table. This proves both identities.
\end{proof}

Designolle and Farkas prove conditional moment identities up to length
five and note that at length six repeated cycles such as $ABCABC$ survive
their reductions, so no general statement is made there \cite[App.~E]{DF}.
Identity (M6) shows that, averaged over outcome strings, exactly these
surviving cycles assemble into the anchored three-fold defects. This
concerns full averages and is used here for heredity; it does not give
the conditional sixth moment needed in \cite[App.~E]{DF}.

\begin{theorem}[Low-order heredity of spectral unbiasedness]
\label{thm:spectralheredity}
Every spectral $k$-UM, at arbitrary finite rank and outcome number, is
pairwise mutually unbiased. For $k\geq3$, every three-setting subtuple
is an algebraic and spectral 3-UM.
\end{theorem}
\begin{proof}
From the second moment alone, \cite[Prop.~6]{DF} obtain pairwise trace
unbiasedness, $\operatorname{Tr}(P_{a|x}P_{b|y})=r/n$; (M4) upgrades this
to the operator identities.
The prescribed characteristic polynomial gives
$\operatorname{Tr}S^m=rp_m(k,n)$ for each selected sum and every $m$.
Apply (M4). Every pair defect is nonnegative, so each vanishes.
A finite average of squared norms vanishes only when every summand
vanishes, giving every pairwise compression identity. Its reverse
follows either by exchanging anchors or by equal-rank singular values.
Now apply (M6). All three anchored defects for every triple vanish,
giving precisely the algebraic 3-UM identities. The spectral triple
conclusion follows from the equivalence \cite[Thm.~3(5)]{DF}.
\end{proof}

\begin{corollary}[Three-outcome spectral no-go at every order]
\label{cor:spectralthree}
There is no spectral $k$-UM with three outcomes at any positive finite
rank for any $k\geq3$.
\end{corollary}
\begin{proof}
Theorem~\ref{thm:spectralheredity} would produce a three-outcome 3-UM,
contradicting the arbitrary-rank three-outcome no-go theorem \cite[Thm.~5]{DF}.
For $k=3$ this is already in \cite{DF}; the new cases are $k\geq4$.
\end{proof}

\section{Quaternionic rigidity at rank two}
Consider an arbitrary strong triple. The standard
unitary-block normal form \cite{FKN} represents every rank-two MUM pair
with computational $P$ by frames whose row blocks are $U_{ab}/2$,
where the $U_{ab}$ are arbitrary complex $2\times2$ unitaries satisfying
\[
 \sum_a U_{ab}^*U_{ab'}=4\delta_{bb'}\id_2,
 \qquad
 \sum_b U_{ab}U_{a'b}^*=4\delta_{aa'}\id_2.
\]
Row and column frame gauges dephase $U_{a0}=U_{0b}=\id_2$.
Explicitly, left-multiply row $a$ by $g_a=U_{a0}^*$ and right-multiply
column $b$ by $h_b=(g_0U_{0b})^*$. These transformations change the
coordinate frames and the outcome frames, respectively, and preserve the
measurements up to a common coordinate change.
Write the row blocks of the third-setting frames as $W_{ac}/2$, with
$W_{ac}\in U(2)$ satisfying the analogous two orthogonality identities.
After the same row transformation, the column frame gauges
$k_c=(g_0W_{0c})^*$ of the third setting set $W_{0c}=\id_2$ without
changing $P,Q$. Thus $U'_{ab}=g_aU_{ab}h_b$,
$W'_{ac}=g_aW_{ac}k_c$ and $S'_{bc}=h_b^*S_{bc}k_c$;
in particular $Z'_{abc}=g_aZ_{abc}g_a^*$.
Unbiasedness of the second and third settings is equivalent to unitarity of
\[
 S_{bc}=\frac12\sum_a U_{ab}^*W_{ac}.
\]
No quaternionic restriction has been imposed on these matrices.

\begin{lemma}[Triangle normal form]\label{lem:triangle}
For pairwise MUM triples in this normal form, the three-fold identities
are equivalent to
\[
 Z_{abc}+Z_{abc}^*=\id_2,
 \qquad Z_{abc}=U_{ab}S_{bc}W_{ac}^*.
\]
If all 64 unitary triangle matrices are nonscalar, the dephased
$U,W,S$ all lie in $SU(2)$ and give a simultaneous quaternionic
realization. 
\end{lemma}
\begin{proof}
The $P_a$ block of $P_aQ_bR_cP_a$ is $Z_{abc}/8$, proving the
first anchored equivalence. Compression to the $Q_b$ and $R_c$ frames
gives unitary conjugates of the same identity or its adjoint, proving
equivalence of the other two anchors.

Unitarity and $Z+Z^*=\id_2$ give $Z^2-Z+\id_2=0$, so its eigenvalues
belong to $\{e^{i\pi/3},e^{-i\pi/3}\}$. A nonscalar $2\times2$
unitary must have distinct eigenvalues here, hence determinant one.
Dephasing gives $Z_{0bc}=S_{bc}$, so $\det S_{bc}=1$.
Then $Z_{a0c}=S_{0c}W_{ac}^*$ forces $\det W_{ac}=1$, and the
general determinant identity forces $\det U_{ab}=1$. The standard
$SU(2)$ realization is the unit-quaternion group. The exact verifier
\texttt{verify\_triangle\_normal\_form.py} checks the constants of this
lemma on two explicit triples (all 64 triangles each): unitary
normalization, Hermitian part and determinant.
\end{proof}

We now remove the nonscalarity hypothesis for four outcomes and rank two.

\begin{lemma}[Balanced pair-phase arrays]\label{lem:phasearray}
Over $\mathbb F_3$, suppose an array with four values for each index has
the form $z_{abc}=u_{ab}+s_{bc}-w_{ac}$ and every line sum is zero.
Then $z_{abc}=0$ everywhere.
\end{lemma}
\begin{proof}
Put $U_b=\sum_a u_{ab}$, $W_c=\sum_a w_{ac}$,
$A_a=\sum_b u_{ab}$ and $U=\sum_b U_b$.
Since $4=1$ in $\mathbb F_3$, summing in $a$ gives
$s_{bc}=W_c-U_b$. Summing in $b$ then gives
$w_{ac}=A_a+W_c-U$.
Hence $z_{abc}=u_{ab}-U_b-A_a+U$ is independent of $c$.
Its sum over $c$ is therefore $4z_{abc}=z_{abc}$, which vanishes by
hypothesis.
\end{proof}

\begin{theorem}[Scalar exclusion and simultaneous quaternionic rigidity]
\label{thm:quaternionrigidity}
Every four-outcome rank-two algebraic 3-UM on $\mathbb C^8$ has no scalar
triangle matrices. In dephased frame coordinates all $U,W,S$ lie in
$SU(2)$, giving three quaternionic orthonormal bases on $\mathbb H^4$.
Equivalently the three PVMs admit a common antiunitary $\Theta$ with
$\Theta^2=-\id_8$ commuting with every outcome projector.
\end{theorem}
\begin{proof}
By Lemma~\ref{lem:triangle}, a triangle $Z$ is unitary with eigenvalues
in $\{e^{i\pi/3},e^{-i\pi/3}\}$. Put $\omega=e^{2\pi i/3}$.
There is a uniquely determined integer $z_{abc}\in\{-1,0,1\}$ such that
\[
 \det Z_{abc}=\omega^{z_{abc}},\qquad
 \operatorname{tr}Z_{abc}=1+i\sqrt3\,z_{abc}.
\]
Here $z=0$ is the nonscalar case; $z=\pm1$ is scalar.

Completeness of the three frame arrays gives
\[
 \sum_a\operatorname{tr}Z_{abc}=4,\qquad
 \sum_b\operatorname{tr}Z_{abc}=4,\qquad
 \sum_c\operatorname{tr}Z_{abc}=4.
\]
For the first equality, cycle the trace and use
$\sum_a W_{ac}^*U_{ab}=2S_{bc}^*$.
For the second, use $\sum_b U_{ab}S_{bc}=2W_{ac}$, and for the third
use $\sum_c S_{bc}W_{ac}^*=2U_{ab}^*$.
These follow directly from the definition of $S$ and the two frame
orthogonality identities. Taking imaginary parts proves that every line
sum of the integer array $z$ is zero.

Dephasing $U_{0b}=U_{a0}=W_{0c}=\id_2$ gives
$\det S_{bc}=\det Z_{0bc}$ and
$\det W_{ac}=\det S_{0c}/\det Z_{a0c}$.
The determinant of any general triangle then shows that all
$\det U_{ab}$ are also powers of $\omega$. Choose exponents
$u_{ab},s_{bc},w_{ac}\in\mathbb F_3$ for these block determinants.
Their multiplicative relation is exactly
$z_{abc}=u_{ab}+s_{bc}-w_{ac}$ modulo three.
Lemma~\ref{lem:phasearray} forces $z=0$ modulo three; since the integer
values are $-1,0,1$, every one is actually zero. Thus scalar triangles
are impossible. Lemma~\ref{lem:triangle} now gives $U,W,S\in SU(2)$.

The standard identification of $SU(2)$ with unit quaternions makes the
frame arrays quaternionic orthonormal bases. In these coordinates
$\Theta=(\id_4\otimes i\sigma_y)\mathcal K$, with $\mathcal K$ complex
conjugation, has square $-\id_8$ and commutes with every quaternionic
block and hence every outcome projector. Conjugating coordinates back
proves the invariant antiunitary statement.
\end{proof}

The determinant-phase calculation is also independently checked by exact
linear algebra over $\mathbb F_3$: the pair-phase array map and its line-sum
map both have rank 37. See
\texttt{verify\_scalar\_triangle\_exclusion.py}. The trace and frame steps
above are analytic and do not rely on searching for scalar examples.
This theorem resolves the scalar regime and gives a global quaternionic
reduction. It does not classify the resulting coupled quaternionic
Hadamard triples or show that an arbitrary first pair is real-Hadamard.

\section{Equivalence of the four-setting definitions}
The following reductions concern four rank-two settings whose every
three-setting subtuple is a strong triple. This is a necessary condition
for an algebraic 4-UM by heredity, and now also for a spectral 4-UM by
Theorem~\ref{thm:spectralheredity}.

\begin{corollary}[A common quaternionic structure for the quartet]
\label{cor:quartetquaternion}
If every three-setting subtuple of four four-outcome rank-two PVMs is an
algebraic 3-UM, all sixteen projectors commute with one antiunitary
$\Theta$ satisfying $\Theta^2=-\id_8$.
\end{corollary}
\begin{proof}
Fix and dephase the first pair once. Apply
Theorem~\ref{thm:quaternionrigidity} to its completion by the third
setting, and then by the fourth. The row gauges for the first pair
are the same in both applications. Only the individual outcome-frame
gauges of the completing settings differ; these do not change their
projectors. Both applications therefore put all frame blocks in the
same copy of $SU(2)$ and yield the same standard antiunitary.
\end{proof}

\begin{theorem}[Conditional algebraic--spectral equivalence at rank two]
\label{thm:conditionalfour}
Under the hypothesis of Corollary~\ref{cor:quartetquaternion}, the quartet
is an algebraic 4-UM if and only if every sum of four selected outcomes
has complex characteristic polynomial
\[
 \left(t^4-4t^3+\frac92t^2-\frac32t+\frac3{32}\right)^2.
\]
\end{theorem}
\begin{proof}
Use the common quaternionic coordinates. Each selected outcome is a
rank-one quaternionic projection $A_i$ on $\mathbb H^4$.
Write $\tau$ for the real quaternionic trace, equivalently half the
complex trace, and $B=\sum_{i=1}^4 A_i$. Pairwise and triple relations give
\[
 \tau(A_i)=1,\quad \tau(A_iA_j)=\frac14,\quad
 \tau(A_iA_jA_k)=\frac1{16}\quad(i,j,k\text{ distinct}).
\]
For the last equality, compress the triple identity by $A_i$ and use
real-trace cyclicity and reversal. Expanding powers gives
\[
 \tau(B)=4,\qquad \tau(B^2)=7,\qquad
 \tau(B^3)=\frac{29}2.
\]
In the fourth power, words with multiplicities $4$, $3+1$, $2+2$,
and $2+1+1$ contribute respectively $4$, $12$, $27/4$, and $9$.
For $2+2$, use $A_iA_jA_i=A_i/4$; for $2+1+1$, cyclicity either
removes an adjacent repeated projector or permits the same pairwise
compression. The remaining contribution is
\[
 4\tau\left(\sum_{\pi\in S_3}
 A_1A_{\pi(2)}A_{\pi(3)}A_{\pi(4)}A_1\right).
\]
The four-fold identity makes this $3/8$, hence
$\tau(B^4)=257/8$. Newton's identities give the displayed quartic
for the four quaternionic eigenvalues. Each occurs twice in the complex
representation, proving the forward direction.

Conversely, the displayed characteristic polynomial fixes
$\tau(B^4)=257/8$. The same word expansion fixes the real trace of
the anchored six-word sum to $3/32$. That sum is Hermitian, commutes
with the common quaternionic structure, and is supported on a rank-one
quaternionic subspace. It is therefore a real scalar times $A_1$.
Its trace proves that scalar is $3/32$. Real-trace cyclicity shows the
same six-cycle trace for every anchor, proving all four-fold identities.
The lower-order identities already hold by hypothesis.
\end{proof}

\begin{corollary}[Equivalence for four outcomes at rank two]
\label{cor:unconditionalfour}
For four four-outcome rank-two PVMs on $\mathbb C^8$, algebraic and
spectral 4-UMs coincide. In particular, every rank-two strong triple
covered by our algebraic nonextension theorems also has no spectral
four-fold extension.
\end{corollary}
\begin{proof}
Algebraic 4-UMs have strong subtriples by heredity. Spectral 4-UMs do
so by Theorem~\ref{thm:spectralheredity}. In either case
Theorem~\ref{thm:conditionalfour} applies. Nonextension follows by
the equivalence. This argument does not extend the equivalence to
higher rank, where quaternionic rank-one compression is unavailable.
By Corollary~\ref{cor:globalnonexistence} both classes turn out to be
empty; the role of this corollary is to transfer the algebraic
obstruction to the spectral definition.
\end{proof}

The equivalence alone does not decide existence.
Theorem~\ref{thm:globalcommuting} below establishes global commuting-pair
necessity, and Corollary~\ref{cor:globalnonexistence} excludes four-setting
tuples in this outcome and rank regime.

\begin{corollary}[A two-moment characterization]
\label{cor:twomoment}
A quartet of four-outcome rank-two PVMs is an algebraic and spectral
4-UM if and only if
\[
 \operatorname{Tr}S_{\vec a}^4=\frac{257}{4}
 \quad\text{for every outcome string},\qquad
 \mathbb E\operatorname{Tr}S_{\vec a}^6=\frac{5441}{16}.
\]
\end{corollary}
\begin{proof}
The target spectrum gives both moments by Newton's identities.
Conversely, the first condition gives equality in (M4), so every pair
is mutually unbiased. The second gives equality in (M6), so every
triple is strong. Common quaternionic rigidity then applies, and the
pointwise fourth moment supplies the anchored four-fold scalar in the
proof of Theorem~\ref{thm:conditionalfour}. No assumption about higher
spectral moments is needed in this converse.
\end{proof}

\section{The commuting-pair extension obstruction}
\begin{theorem}[Uniform commuting-pair nonextension]
\label{thm:commutingnonextension}
Every four-outcome rank-two strong triple containing a pair with
commuting canonical blocks forbids an algebraic or spectral fourth
setting.
\end{theorem}
\begin{proof}[Outline of the proof]
Quaternionic rigidity and the order-four complex Hadamard classification
put a commuting pair in a one-parameter normal form. For every nonreal
phase, the strong-triple equations classify third outcomes into eight
signed circles; orthogonality reduces complete third measurements to one
normal form. The fourth-setting moment constraints force its centered
outcomes $D=T-\id/4$ to anticommute with a fixed involution, giving
$\operatorname{Tr}D^3=0$. A rank-two projector in dimension eight instead
has $\operatorname{Tr}(T-\id/4)^3=3/4$, a contradiction.
An explicit positive Gram factorization makes this argument uniform over
all parameters. The real-phase case follows from the fixed-real-pair
classification and obstruction. Appendix~\ref{app:commutingproof} gives
the complete classification and algebra; Appendix~\ref{app:fixedpair}
handles the real-phase case. Spectral nonextension follows from
Corollary~\ref{cor:unconditionalfour}.
For a commuting pair the antiunitary commuting with both settings is not
unique; throughout we use the standard one in the shared dephased
coordinates of the first pair, as in Corollary~\ref{cor:quartetquaternion}.
\end{proof}

\section{Global obstruction at rank two}
\subsection{Normal form of a noncommuting first pair}
For a Hermitian involution with zero diagonal in a four-outcome basis,
its \emph{support graph} has an edge $ab$ when $X_{ab}\ne0$.
Call the support full when all six edges are present.
\begin{lemma}[Coarse-support dichotomy]\label{lem:coarsesupport}
For any quaternionic four-outcome MUM pair $P,Q$, consider the three
balanced involutions
$Q_0+Q_1-Q_2-Q_3$, $Q_0-Q_1+Q_2-Q_3$ and
$Q_0-Q_1-Q_2+Q_3$ in the $P$ basis.
If the pair's canonical blocks commute, all three supports are non-full.
If they do not commute, at most one support is non-full.
In particular, if any one of the three supports is a matching, the
canonical blocks commute.
\end{lemma}
\begin{proof}
Two distinct vertices cannot have exactly one common neighbor, since
then an off-diagonal entry of $X^2$ would be a single nonzero product.
Every vertex has positive degree. The only possible graphs on four
vertices are therefore a matching, a four-cycle or the full graph.

A matching involution can be row-gauged and permuted to
$X=\bigl(\begin{smallmatrix}0&1&0&0\\1&0&0&0\\
0&0&0&1\\0&0&1&0\end{smallmatrix}\bigr)$.
For either other balanced involution $Y$, commutation with $X$ and the
zero diagonals of $Y,XY$ give
\[
 Y=\begin{pmatrix}0&0&a&b\\0&0&b&a\\
 \bar a&\bar b&0&0\\\bar b&\bar a&0&0\end{pmatrix},
 \qquad |a|^2+|b|^2=1,\quad\operatorname{Re}(a\bar b)=0.
\]
A further gauge preserving $X$ makes one nonzero entry real and the
other purely imaginary. All three involutions, and hence the four
projectors, lie in a common complex subfield. The canonical pair commutes.

For a four-cycle, a tree gauge makes three edges real, and $X^2=I$
forces the remaining cycle product to be negative real. After relabeling,
\[
 X=\begin{pmatrix}0&0&r&s\\0&0&s&-r\\r&s&0&0\\s&-r&0&0\end{pmatrix},
 \qquad rs\ne0,\quad r^2+s^2=1.
\]
The same commutation and diagonal conditions give, exhaustively,
\[
 Y=\begin{pmatrix}
 0&a&sf&-rf\\-a&0&-rf&-sf\\
 -sf&rf&0&-a\\rf&sf&a&0
 \end{pmatrix},\qquad
 \bar a=-a,\quad\bar f=-f.
\]
Here $Y^2=I$ gives $|a|^2+|f|^2=1$ and
$\operatorname{Re}(a\bar f)=0$. If $Y$ is also non-full, then $a=0$
or $f=0$, making $XY$ or $Y$, respectively, a matching involution.
Thus two non-full supports force canonical commutativity.

Finally, a full-support involution can be row-gauged to
\begin{equation}\label{eq:fullcoarsenormal}
 X=\begin{pmatrix}
 0&r&s&t\\r&0&t\mathbf i&-s\mathbf i\\
 s&-t\mathbf i&0&r\mathbf i\\
 t&s\mathbf i&-r\mathbf i&0
 \end{pmatrix},\quad r,s,t>0,\quad r^2+s^2+t^2=1.
\end{equation}
To see this, make the first-row edges positive real; the off-diagonal
square equations force the three remaining edges to be scalar multiples
of one imaginary quaternion, and the off-diagonal equations among the
indices $1,2,3$ give their magnitudes $t,s,r$ (for example
$|X_{23}|=r$); the diagonal equations give $r^2+s^2+t^2=1$.
A common conjugation sets the direction to $\mathbf i$.
The only imaginary block-diagonal complex structure commuting with this
$X$ is, up to sign, $J=\mathbf i I$.
On either $e_\pm$, put $T_a=e_\pm P_a e_\pm-e_\pm/4$. A direct
computation gives the real Gram matrix $\tau(T_aT_b)=1/8$ for $a=b$ and
$|X_{ab}|^2/4-1/8$ for $a\neq b$. Its eigenvectors are the constant vector
and the three Walsh sign vectors, with eigenvalues $0,r^2/2,s^2/2,t^2/2$.
Only the moduli $|X_{ab}|$ enter, so for $X=\Gamma(n)$ (notation of
Section~\ref{sec:globalexhaustion}) with any unit vector $n$, including
vectors with negative or zero coordinates, the eigenvalues are $0$ and
$n_i^2/2$. They span the three-dimensional traceless
Hermitian subspace commuting with $J$. A traceless supported matrix
with zero $P$ diagonal therefore anticommutes with $J$.
If the canonical pair $P,Q$ commuted, its blocks would lie in one complex
subfield $\mathbb C_\mu$ in some gauge, and $\mu I$ would be a complex
structure commuting with both settings. That gauge and the full-support
gauge both make $P$ computational, so they differ by a monomial
quaternionic unitary, which sends $\mu I$ to a diagonal of imaginary
units. Hence this complex structure would have to equal $J$ up to sign. But the
nonzero differences of $Q$ outcomes supported on $e_\pm$ would then
both commute and anticommute with $J$, a contradiction.
This excludes full support for a commuting pair and finishes the dichotomy.
\end{proof}

\begin{corollary}[One projector chart for every noncommuting first pair]
\label{cor:unifiedpairchart}
Every noncommuting quaternionic four-outcome MUM pair is equivalent,
up to coordinate gauge and outcome relabelling, to the following
three-parameter family. Choose real variables with
\[
 r^2+h^2=c^2+d^2=a^2+b^2=1,\qquad r,h,c,d>0.
\]
Let $V$ be the full-support matrix \eqref{eq:fullcoarsenormal} with
$s=hc$, $t=hd$, and put $E_\sigma=(I+\sigma V)/2$. Define complex
skew-symmetric matrices by their upper-triangular entries:
\[
\begin{aligned}
 (G_\sigma)_{01}&=h/2,&
 (G_\sigma)_{02}&=(-rc-\sigma i d)/2,&
 (G_\sigma)_{03}&=(-rd+\sigma i c)/2,\\
 (G_\sigma)_{12}&=(-\sigma c-i rd)/2,&
 (G_\sigma)_{13}&=(-\sigma d+i rc)/2,&
 (G_\sigma)_{23}&=i h/2.
\end{aligned}
\]
Here $i$ denotes the same complex subfield as $\mathbf i$.
With $A_+=(a+i b)G_+\mathbf j$ and
$A_-=(a-i b)G_-\mathbf j$, the second
setting is
\[
 Q_0=(E_++A_+)/2,\quad Q_1=(E_+-A_+)/2,\quad
 Q_2=(E_-+A_-)/2,\quad Q_3=(E_--A_-)/2.
\]
All entries are polynomial in the six circle coordinates. The
parametrization need not be unique. It is a parametrization of order-four
quaternionic Hadamard matrices with a noncommuting canonical form;
compare the classification of \cite[Thm.~4.6]{CD} and the use of these
families in \cite[Sec.~2.5]{FKN}.
\end{corollary}
\begin{proof}
A noncommuting pair has at least two full balanced coarse supports by
Lemma~\ref{lem:coarsesupport}. Select one and use its full-support gauge.
The differences of the two outcomes on each $E_\sigma$ are supported,
zero-$P$-diagonal Hermitian involutions. By the last paragraph of the
proof of Lemma~\ref{lem:coarsesupport}, each anticommutes with $J$, so it
has the form $D\mathbf j$ with $D$ complex skew-symmetric. Support on
$E_\sigma$ places $D$ in the exterior square of the two-dimensional range
of $E_\sigma$, a one-dimensional complex space spanned by $G_\sigma$.
Hence each difference is $\lambda_\sigma G_\sigma\mathbf j$, and the
involution condition gives $|\lambda_\sigma|=1$.
Conjugation by a common scalar complex phase preserves $P,V$ and rotates
both circles by the same angle. Remove their mean angle, giving opposite
phases $a+i b$ and $a-i b$. Every relative unit phase has a square root,
so this symmetric convention covers every pair. This proves coverage.
Conversely, the displayed matrices satisfy $A_\sigma^2=E_\sigma$ and
$VA_\sigma=\sigma A_\sigma$. They give four quaternionic rank-one
orthogonal projectors, each with $P$ diagonal $1/4$, hence a MUM pair.
Its full balanced support excludes canonical commutativity.
The symbolic verifier checks these projector and MUM identities modulo
the three circle equations.
\end{proof}

\subsection{Full proof of global commuting-pair necessity}\label{sec:globalexhaustion}
Write $\Gamma(m)$ for the matrix in \eqref{eq:fullcoarsenormal}, with
$m=(r,hc,hd)$, and use $\Gamma$ for the real-linear extension of this
formula to all of $\mathbb R^3$. It satisfies
$\Gamma(a)\Gamma(b)+\Gamma(b)\Gamma(a)=2(a\cdot b)I$, which gives
$E_n\Gamma(m)E_n=(n\cdot m)E_n$ below. Use the symmetric convention of
Corollary~\ref{cor:unifiedpairchart}:
\[
 A_+=(a+ib)G_+\mathbf j,\qquad A_-=(a-ib)G_-\mathbf j,
 \qquad a^2+b^2=1.
\]
Throughout this section $P$ and $Q$ are the fixed first two settings;
$T$ denotes one outcome of a prospective third setting. Define
\[
 u=(ah,-arc+bd,-ard-bc),\quad
 v=(bh,-brc-ad,-brd+ac),\quad w=-m.
\]
These vectors form an oriented orthonormal basis, $u\times v=w$.
The first-row entries of $2(a+ib)G_+$ are $u+iv$; the opposite entries
are $i(u_0+iv_0),-i(u_1+iv_1),i(u_2+iv_2)$.
For $G_-$ they are the complex conjugates of the first-row entries,
and the negatives of the conjugates of the opposite entries.


\begin{lemma}[Normalization of the single-zero stratum]
\label{lem:stratumsymmetry}
In the symmetric chart above, if exactly one coordinate among $u,v$
vanishes, allowed coordinate gauges and outcome relabellings preserve
the selected full coarse-graining and put the chart in the form
\[
 v_0=0,\qquad a=1,\quad b=0,\qquad
 u=(h,-rc,-rd),\quad v=(0,-d,c).
\]
The transformed parameters still satisfy $r,h,c,d>0$ and both circle
identities. These operations preserve every strong-triple identity.
\end{lemma}
\begin{proof}
Put $B=A_++A_-$ and $C=A_+-A_-$, the other two balanced
coarse-grainings of $Q$. Their first-row entries are respectively
$u_i\mathbf j$ and $v_i\mathbf k$; their opposite entries are
$\epsilon_i u_i\mathbf k$ and $-\epsilon_i v_i\mathbf j$,
where $\epsilon=(1,-1,1)$ refers to edges $23,13,12$.
A permutation $\rho$ of the last three computational coordinates
permutes the three coordinates of $m$. For an odd permutation,
simultaneous conjugation by $\mathbf j$ restores the displayed
$\Gamma(m)$ convention, since it reverses $\mathbf i$ and
$\mathbf k$ while fixing $\mathbf j$.
Consequently the transformed vectors are
$u'=\rho u$, $v'=(\operatorname{sgn}\rho)\rho v$, and
$m'=\rho m$, with the same positive first-row gauge.
The transformed pair is again exactly in the symmetric chart: the
supported skew part on each $E'_\sigma$ is one complex dimension, so
its first row determines it, and since the $B$ entries are multiples of
$\mathbf j$ and the $C$ entries multiples of $\mathbf k$ the two phases
remain complex conjugate. Orientation is kept because
$\rho u\times(\operatorname{sgn}\rho)\rho v=-\rho m$.
Reparametrize it by $r=m'_0$, $h=\sqrt{(m'_1)^2+(m'_2)^2}$,
$c=m'_1/h$, $d=m'_2/h$; all four remain positive.
Choose $\rho$ to move the zero to index zero.

If that zero belongs to $u$, interchange $Q_2,Q_3$, so
$A_-\mapsto-A_-$ and $B,C$ are exchanged. Simultaneous conjugation
by $\exp(-\mathbf i\pi/4)$ then gives
$(u,v)\mapsto(v,-u)$: it sends $\mathbf j$ to $-\mathbf k$
and $\mathbf k$ to $\mathbf j$, and fixes $\Gamma(m)$.
Thus in either case $v_0=0$. In the symmetric chart $v_0=bh$ and
$h>0$, hence $b=0$ and $a=\pm1$.
If needed, interchange the two outcomes on both $E_+$ and $E_-$;
this negates both $A_\pm$, replacing $(a,b)$ by $(-a,-b)$.
We may therefore take $a=1$, which gives the stated vectors.
All operations are simultaneous quaternionic unitary gauges or
outcome permutations. They preserve the measurement equations and
the selected coarse-graining, up to its explicitly recorded coordinate
permutation and gauge.
\end{proof}

\begin{theorem}[Global commuting-pair necessity]
\label{thm:globalcommuting}
Every four-outcome rank-two strong triple contains a pair with commuting
canonical blocks, equivalently a pair that is a direct sum of mutually
unbiased bases.
\end{theorem}
\begin{proof}
Quaternionic rigidity reduces to quaternionic rank-one projectors. If
any pair commutes, there is nothing to prove. Otherwise fix its chart
$P,Q$ above and classify every possible third outcome $T$.
Write $T=C+D\mathbf j$, where $C$ is complex Hermitian and $D$ complex
skew-symmetric. The $P$ diagonal of $C$ is $1/4$. Put
\[
 \xi=(D_{01},D_{02},D_{03})=x+iy.
\]
By Proposition~\ref{prop:globalkernels}, the strong equations are
equivalent to the real-trace moment conditions for $T-I/4$.
Substituting the displayed first-pair chart and $T=C+D\mathbf j$,
these split into the complex conditions
\begin{equation}\label{eq:globalClinear}
 \operatorname{diag}(C\Gamma(m)+\Gamma(m)C)=0
\end{equation}
and real row-contraction conditions against $(a\pm ib)G_\pm$.
To make the linear reduction explicit, let the opposite edges to
$01,02,03$ be $23,13,12$, and put $\epsilon=(1,-1,1)$.
Write their $D$ entries as $\zeta_i=s_i+it_i$. The first-row and
opposite-edge contraction weights against the two phased $G$ matrices
are, respectively,
\[
 \tfrac12(u_i x_i+v_i y_i),\quad
 \tfrac12\epsilon_i(u_i t_i-v_i s_i),\qquad
 \tfrac12(u_i x_i-v_i y_i),\quad
 \tfrac12\epsilon_i(u_i t_i+v_i s_i).
\]
For symmetric real edge weights on four vertices, zero row sums are
equivalent to equality of opposite weights and a zero sum of the three
first-row weights: signed combinations of the four row equations give
the opposite equalities, and their common remaining equation is that sum.
Adding and subtracting the two phased systems therefore gives
\[
 u_i(x_i-\epsilon_i t_i)=0,\qquad
 v_i(y_i+\epsilon_i s_i)=0,\qquad
 \sum_i u_i x_i=\sum_i v_i y_i=0.
\]
Consequently, if $u_i v_i\ne0$ for every $i$, these conditions are exactly
\begin{equation}\label{eq:globalDlinear}
 D_{23}=i\xi_0,\quad D_{13}=-i\xi_1,\quad D_{12}=i\xi_2,
 \qquad x\cdot u=y\cdot v=0.
\end{equation}
At an index where $v_i=0$ the real part of the corresponding opposite
entry is free; the analogous statement holds when $u_i=0$.
There is at most one zero among the six coordinates of $u,v$.
Indeed, two zeros in one vector contradict its orthogonality to the
strictly nonzero vector $w$. Zeros in different vectors at different
indices force a second zero by $u\cdot v=0$; at the same index they
force $w$ to be a coordinate vector. This exhausts the possible strata.

\emph{Complex outcomes.} If $D=0$, choose a normalized frame
$(1,e^{i\alpha},e^{i\beta},e^{i\gamma})/2$ for $T$.
Equation~\eqref{eq:globalClinear} gives
\[
 \cos\alpha=\sin(\beta-\gamma),\quad
 \cos\beta=\sin(\gamma-\alpha),\quad
 \cos\gamma=\sin(\alpha-\beta).
\]
For the vectors of cosines and sines this says $A=B\times A$.
Its scalar product with $A$ forces $A=0$. Thus the eight frames are
\begin{equation}\label{eq:globaleightcomplex}
 (1,i\varepsilon_0,i\varepsilon_1,i\varepsilon_2)/2,
 \qquad \varepsilon_i\in\{1,-1\}.
\end{equation}
They split into two orthonormal settings according to the product of
the three signs. Both commute canonically with $P$.

\emph{Nonzero $D$ without a zero coordinate.}
Since $T$ has quaternionic rank one, $D$ has complex rank two and
vanishing Pfaffian. Under \eqref{eq:globalDlinear} its Pfaffian is
$i\xi\cdot\xi$, so $x\cdot y=0$ and $|x|=|y|>0$.
Write
\[
 x=\alpha v+\beta w,\qquad y=\gamma u+\delta w.
\]
These relations give $\beta\delta=0$ and
$\alpha^2+\beta^2=\gamma^2+\delta^2$.
An entrywise multiplication gives
\[
 DD^*=N I-2\Gamma(x\times y)=2N E_n,\qquad
 N=|x|^2+|y|^2,\quad E_n=(I+\Gamma(n))/2,
 \quad n=-\frac{x\times y}{|x|^2}.
\]
For clarity, write a normalized quaternionic outcome frame as
$Z+W\mathbf j$, with complex columns $Z,W$. Then
$C=ZZ^*+WW^*$ and $D=WZ^{\mathsf T}-ZW^{\mathsf T}$.
Since $D\ne0$, the columns are independent, so both matrices have the
same two-dimensional complex range, namely $E_n$.
Because $E_n\Gamma(m)E_n=(n\cdot m)E_n$ and
$\operatorname{tr}C=1$, summing \eqref{eq:globalClinear} gives
$n\cdot m=0$, equivalently $\alpha\gamma=0$.
Together these equations force $n=\pm u$ or $n=\pm v$ and
\[
 \xi=t(v-isw)\quad\text{or}\quad \xi=t(w-isu),
 \qquad s\in\{1,-1\},\quad t\in\mathbb R\setminus\{0\}.
\]
All coordinates of $n$ are nonzero. The compressed-diagonal span in
Lemma~\ref{lem:coarsesupport}, applied to $E_n$, forces $C=E_n/2$.
The complex part of $T^2=T$ then gives $DD^*=E_n/4$, hence $t=\pm1/4$.
There are exactly four outcomes of each type.

\emph{The single-zero stratum.}
By Lemma~\ref{lem:stratumsymmetry}, we may take $a=1,b=0$
without omitting any single-zero case or changing the strong identities.
Thus, exhaustively,
\[
 u=(h,-rc,-rd),\quad v=(0,-d,c),\quad w=(-r,-hc,-hd),
\]
and the linear conditions replace \eqref{eq:globalDlinear} only by
\[
 D_{23}=i\xi_0+\kappa,\qquad \kappa\in\mathbb R.
\]
The remaining two opposite entries and the orthogonality conditions
are unchanged. The Pfaffian equations become
\begin{equation}\label{eq:globalexceptioncone}
 \beta(2\delta+r\kappa)=0,\qquad
 \alpha^2+\beta^2-\gamma^2-\delta^2
       +\kappa(h\gamma-r\delta)=0.
\end{equation}

If $\kappa=0$, the preceding argument again gives $n=\pm u,\pm v$.
Only $n=\pm v$ has a zero coordinate. The supported complex Hermitian
matrices with zero diagonal in $E_{sv}$ form the real line generated by
\[
 H_s=\begin{pmatrix}
 0&-1&-isc&-isd\\-1&0&sd&-sc\\
 isc&sd&0&i\\isd&-sc&-i&0
 \end{pmatrix}.
\]
The Gram computation in Lemma~\ref{lem:coarsesupport}, specialized
to $n=sv$ with one zero coordinate, has eigenvalues
$0,0,c^2/2,d^2/2$, of which exactly two are nonzero.
It therefore leaves one traceless complex Hermitian direction on the
rank-two space $E_{sv}$; the displayed $H_s$ spans that direction.
The verifier independently confirms this one-dimensional kernel by an
exact rank-eleven minor.
Its Jordan diagonal with $\Gamma(m)$ is $r(-2,-2,2,2)$.
Thus \eqref{eq:globalClinear} kills this extra line, since $r>0$,
and the same eight noncomplex outcomes remain.

Suppose now $\kappa\ne0$. Dephase the quaternionic outcome frame as
$(1,z_0-4\xi_0\mathbf j,z_1-4\xi_1\mathbf j,
 z_2-4\xi_2\mathbf j)/2$, where $z=A+iB$ and $A,B$ are real.
The interior $D$ entries and \eqref{eq:globalClinear} give
\begin{equation}\label{eq:globalframeequations}
 \xi\times z=-i\xi-\kappa e_0,\qquad
 A+A\times B+16x\times y=0,\qquad m\cdot A=0.
\end{equation}
The three coordinate norms are
\begin{equation}\label{eq:globalframenorms}
 A_i^2+B_i^2+16(x_i^2+y_i^2)=1.
\end{equation}
Write $A=p u+q v$ and $B=t u+\ell v+n w$.
The real and imaginary parts of the cross equation are, respectively,
\begin{align*}
 -\beta q+\delta\ell&=\gamma-\kappa h,&
 \beta p+\gamma n-\delta t&=0,&
 -\alpha p-\gamma\ell&=\delta+\kappa r,\\
 \alpha n-\beta\ell-\delta q&=0,&
 \beta t+\delta p&=-\alpha,&
 -\alpha t+\gamma q&=-\beta.
\end{align*}
The remaining strong equations are
\begin{equation}\label{eq:globalstrongcoordinates}
 p+qn+16\alpha\delta=0,\quad
 q-pn+16\beta\gamma=0,\quad
 p\ell-qt-16\alpha\gamma=0.
\end{equation}
There are only two branches in \eqref{eq:globalexceptioncone}.

If $\beta=0$ and $\alpha\ne0$, then $\delta\ne0$, and the linear
system gives
$p=-\alpha/\delta$, $q=\alpha n/\delta$,
$t=\gamma n/\delta$, $\ell=(\gamma-\kappa h)/\delta$.
Equations~\eqref{eq:globalstrongcoordinates} force
$n=0$, $\delta^2=1/16$ and $\gamma=h\kappa/2$.
The Pfaffian gives
$\alpha^2=\delta^2+\kappa r\delta-\kappa^2h^2/4$.
Substituting in the last two norms yields $1+16\kappa r\delta=1$,
a contradiction. Thus this branch has $\alpha=\beta=0$.

On the other branch $\delta=-r\kappa/2$, put
$g=\gamma-h\kappa/2$ and $L=\beta^2+\delta^2>0$.
The Pfaffian gives
$g^2=\alpha^2+\beta^2+\kappa^2/4$ and
$L+\alpha^2=\gamma(\gamma-h\kappa)$.
In particular neither factor on the right vanishes. The linear system
has the uniform solution
\[
 p=-\frac{\alpha\delta+\beta\gamma n}{L},\quad
 q=\frac{\alpha\delta n-\beta(\gamma-h\kappa)}{L},\quad
 t=\frac{\delta\gamma n-\alpha\beta}{L},\quad
 \ell=\frac{\delta(\gamma-h\kappa)+\alpha\beta n}{L}.
\]
The last strong equation reduces to
\[
 -\frac{\alpha}{L}
 \bigl[\gamma(n^2+16L+1)-h\kappa\bigr]=0.
\]
The bracket can vanish only if $h\kappa/\gamma=n^2+16L+1>1$. But
$h\kappa/\gamma<1$: since $g^2\geq\kappa^2/4$, we have
$|g|\geq|\kappa|/2>h|\kappa|/2$, so $\gamma=g+h\kappa/2$ has the sign of
$g$. If $g\kappa<0$ the ratio is therefore negative; otherwise
$g/\kappa\geq1/2$ and the ratio is at most $2h/(1+h)<1$. Since $L>0$,
this forces $\alpha=0$. Then the first strong equation reads
$-\beta n(2\gamma-h\kappa)/L=0$, and $2\gamma-h\kappa=2g\neq0$. So if
$\beta\ne0$, the first two strong equations give $n=0$ and
$\gamma^2=1/16$. The last two norms become
$1-h\kappa/\gamma=1$, again impossible. Hence $\beta=0$ as well.
All nonzero-$\kappa$ solutions therefore have $x=0$.

Taking a scalar product of the middle equation in
\eqref{eq:globalframeequations} with $A$ now gives $A=0$.
Write $y=(f,ec,ed)$, as required by $y\cdot v=0$.
The cross equation excludes $y=0$. The Pfaffian gives
$\kappa f=f^2+e^2$, so $f\ne0$.
The last two cross equations give
\[
 B_1=(e/f)(cB_0-d),\qquad B_2=(e/f)(dB_0+c).
\]
The first norm is $B_0^2+16f^2=1$; the other two reduce to
$(e^2/f^2)(1-2cdB_0)=1$ and
$(e^2/f^2)(1+2cdB_0)=1$.
They force $B_0=0$, $e=\pm f$, $f=\pm1/4$ and $\kappa=2f$.
The four additional frames are consequently
\begin{equation}\label{eq:globalmixedfour}
 (1,\varepsilon\mathbf k,
 \eta(-d\mathbf i+\varepsilon c\mathbf k),
 \eta(c\mathbf i+\varepsilon d\mathbf k))/2,
 \qquad \varepsilon,\eta\in\{1,-1\}.
\end{equation}
This finishes outcome exhaustion, including all exceptional kernels.

\emph{Completeness and canonical commutativity.}
The sixteen ordinary outcomes split into four sets of four: the two
sign-product classes in \eqref{eq:globaleightcomplex}, and the $u$ and
$v$ types. Distinct outcomes in one class have overlap
$\tau(TT')=0$; outcomes in different ordinary classes have overlap $1/4$.
The exceptional four in \eqref{eq:globalmixedfour} are also mutually
orthogonal. Their overlaps with the two complex classes and the $u,v$
classes are respectively among
$(1\pm c)/4$, $(1\pm d)/4$, $(1\pm h)/4$, $(1\pm r)/4$.
All are strictly positive because $r,h,c,d\in(0,1)$.
Thus the orthogonality graph is four disjoint complete graphs on four
vertices, or five on the exceptional stratum. A complete measurement
must be one whole class; mixed settings are impossible.

For the two complex classes, $P,T$ commute canonically.
For the exceptional class, dephasing the displayed frames places all
entries in $\mathbb R+\mathbb R\mathbf j$, again giving a commuting
pair with $P$. For the remaining two classes their balanced involution
is $\Gamma(u)$ or $\Gamma(v)$. Direct multiplication gives
\[
\begin{gathered}
 \Gamma(u)A_+\Gamma(u)=-A_-,\qquad
 \Gamma(v)A_+\Gamma(v)=A_-,\\
 \Gamma(u)\Gamma(m)=-\Gamma(m)\Gamma(u),
 \qquad\Gamma(v)\Gamma(m)=-\Gamma(m)\Gamma(v).
\end{gathered}
\]
Therefore in the $Q$ basis this involution swaps the two rank-one
outcomes on $E_+$ with those on $E_-$: its support is a matching.
By the matching case of Lemma~\ref{lem:coarsesupport}, $Q,T$ commute canonically.
Every complete setting has the required commuting pair.
\end{proof}

\begin{corollary}[Nonexistence of four-fold unbiased measurements]
\label{cor:globalnonexistence}
There is no algebraic or spectral four-fold unbiased tuple with four
outcomes of rank two on $\mathbb C^8$.
\end{corollary}
\begin{proof}
Any three-setting subtuple is strong by algebraic heredity or
Theorem~\ref{thm:spectralheredity}. It contains a commuting canonical
pair by Theorem~\ref{thm:globalcommuting}, contradicting the uniform
nonextension theorem, Theorem~\ref{thm:commutingnonextension}.
\end{proof}

The classification is global in the first pair. Generic and exceptional
strata refer to the exact vanishing of the displayed chart coordinates;
all noncommuting pairs are included. Chart endpoints at which the chosen
coarse-graining loses an edge are handled by choosing a full
coarse-graining, whose existence follows from the coarse-support
lemma. Commuting first pairs require no additional classification for
Theorem~\ref{thm:globalcommuting} and are covered uniformly by the
nonextension theorem.

The result is specific to rank two. Full algebraic--spectral equivalence
and unrestricted nonexistence at higher rank remain open here: the
quaternionic rank-one reduction is essential to the outcome exhaustion.

\paragraph{Proof verification.}
The main moment argument is analytic; its word counts and Newton
identities are checked by exact rational arithmetic. The quaternionic
rigidity proof is a determinant-phase argument over $\mathbb F_3$.
The extension theorem uses exact polynomial identities and a uniformly
positive Gram factorization. The global classification uses analytic
case exhaustion, with exact symbolic checks of its matrix identities,
exceptional reductions, candidate overlaps and strong moments. Appendix~\ref{app:verification} identifies
the independently executable verifiers and the additional certificates,
which are available at
\url{https://github.com/danielrupawalla/rank-two-obstruction}.

\paragraph{AI assistance.}
This manuscript was developed through human-directed discussions with
OpenAI Codex. The AI assistant contributed proposed proof strategies,
algebraic derivations, verifier code, internal consistency checks, and
writing and editing of the text. Daniel Rupawalla selected the research
questions and directed the work. The author is responsible for
checking and approving the final mathematical claims and manuscript.

\begin{thebibliography}{9}
\bibitem{DF} S. Designolle and M. Farkas,
\emph{$k$-fold unbiased measurements and maximal incompatibility},
arXiv:2609.20728v1 (2026).
\url{https://arxiv.org/abs/2609.20728}.
\bibitem{F} M. Farkas, \emph{$n$-fold unbiased bases: an extension of the MUB condition},
arXiv:1706.04446 (2017). \url{https://arxiv.org/abs/1706.04446}.
\bibitem{FKN} M. Farkas, J. Kaniewski and A. Nayak,
\emph{Mutually Unbiased Measurements, Hadamard Matrices, and Superdense Coding},
IEEE Trans. Inf. Theory \textbf{69}, 3814 (2023);
arXiv:2204.11886. \url{https://arxiv.org/abs/2204.11886}.
\bibitem{CD} O. Chterental and D. \v{Z}. Djokovi\'c,
\emph{On orthostochastic, unistochastic and qustochastic matrices},
Linear Algebra Appl. \textbf{428}, 1178 (2008);
arXiv:0708.4051. \url{https://arxiv.org/abs/0708.4051}.
\bibitem{TZ} W. Tadej and K. \.{Z}yczkowski,
\emph{A concise guide to complex Hadamard matrices},
Open Syst. Inf. Dyn. \textbf{13}, 133--177 (2006).
\url{https://arxiv.org/abs/quant-ph/0512154}.
\bibitem{DS} F. Del Santo, J. Czartowski, K. \.{Z}yczkowski and N. Gisin,
\emph{Iso-entangled bases and joint measurements},
Phys. Rev. Research \textbf{6}, 023085 (2024).
\url{https://doi.org/10.1103/PhysRevResearch.6.023085}.
\end{thebibliography}

\appendix
\section{Proof of the commuting-pair obstruction}\label{app:commutingproof}
We now strengthen the fixed real-Hadamard result to every commuting
canonical pair within a strong triple. The standard pairwise
commutativity criterion is from \cite{FKN}; its necessity somewhere in
an arbitrary strong triple is not assumed.

\begin{lemma}[Commuting first-pair normal form]\label{lem:commutingpair}
A pair within a four-outcome rank-two strong triple, whose dephased
canonical blocks commute, is equivalent to the pair with computational
first setting and normalized second-setting frame array $U(q)/2$, where
\[
 U(q)=\begin{pmatrix}
 1&1&1&1\\1&-1&q&-q\\1&1&-1&-1\\1&-1&-q&q
 \end{pmatrix},\qquad q=c+s\mathbf i,\quad c^2+s^2=1.
\]
\end{lemma}
\begin{proof}
By quaternionic rigidity, the canonical blocks lie in $SU(2)$.
Commuting blocks lie in one complex subfield of the quaternions after
a common internal conjugation. Their scalar complex entries form an
order-four complex Hadamard matrix. The standard classification puts
this in the one-parameter family \cite[Section~5.4]{TZ}. Row and column
phase gauges in that subfield act simultaneously on the conjugate
complex component, so they are allowed quaternionic frame gauges.
Permuting rows and columns gives the displayed form.
\end{proof}

\begin{lemma}[Third-outcome classification for $s\ne0$]
\label{lem:commutingoutcomes}
Normalize a quaternionic third-outcome frame to
$\frac12(1,u_1,u_2,u_3)^T$. Identify imaginary quaternions with
$\mathbb R^3$, put $n=\mathbf i$, and let
$v=\cos t\,\mathbf j+\sin t\,\mathbf k$ and $h=n\times v$.
Its pair and strong triple identities with the preceding pair hold
exactly in these two classes, with all displayed signs independent:
\[
\begin{array}{ll}
\mathrm A_{\varepsilon,\delta}(t):&
 (u_1,u_2,u_3)=(v,\varepsilon n,\delta h),\\[2pt]
\mathrm B_{\eta,\varepsilon}(t):&
 u_1=\frac{\eta s}{2}+\frac{\varepsilon-\eta c}{2}n
       +\frac{\varepsilon+\eta c}{2}h-\frac{\eta s}{2}v,\\
&u_2=v,\\
&u_3=-\frac{\eta s}{2}+\frac{\varepsilon+\eta c}{2}n
       +\frac{\varepsilon-\eta c}{2}h+\frac{\eta s}{2}v.
\end{array}
\]
Each sign belongs to $\{1,-1\}$.
\end{lemma}
\begin{proof}
For column $b$, put $v_a^{(b)}=U_{ab}^*u_a$, with $u_0=1$.
The first-anchor triple identities say that the real Gram matrix of
the four unit quaternions $v_a^{(b)}$ has every row sum equal to one.
Indeed $S_b=\frac12\sum_a v_a^{(b)}$ and
$\operatorname{Re}(S_b\overline{v_a^{(b)}})=1/2$.
These row sums also give $|S_b|=1$, so the pair relation with the
second setting holds and the triangle normal form gives all anchors.

The first two columns of $U(q)$ imply
\[
 u_1=a+x,\quad u_2=v_0,\quad u_3=-a+y,
 \qquad |v_0|=1,\quad
 x\cdot v_0=-a,\quad y\cdot v_0=a,\quad x\cdot y=a^2,
\]
where $x,y,v_0$ are imaginary and $|x|^2=|y|^2=1-a^2$.
Thus $p=(x+y)/\sqrt2$ is a unit vector orthogonal to $v_0$.
With $w=v_0\times p$ there is a signed real $b$ such that
\[
 x-y=-2av_0+\sqrt2 b w,\qquad b^2+4a^2=1.
\]
The last two columns give
\[
 n\cdot p=n\cdot w,\qquad a(n\cdot v_0)=0,\qquad
 2ac+s\bigl(\sqrt2 b(n\cdot w)-2a(n\cdot v_0)\bigr)=0.
\]
These follow by expanding real quaternionic products; the converse
recovers all four Gram row-sum identities.
For clarity, put $\alpha=n\cdot v_0$, $\gamma=n\cdot p$ and
$\delta=n\cdot w$. The four row sums minus one for column 2 are
\[
\begin{split}
 r_0&=2ac+s(-2a\alpha+\sqrt2 b\delta),\\
 r_1&=2ac+\frac{s}{\sqrt2}((1+b)\gamma+(b-1)\delta),\\
 r_2&=2ac+s(2a\alpha+\sqrt2 b\gamma),\\
 r_3&=2ac+\frac{s}{\sqrt2}((b-1)\gamma+(b+1)\delta).
\end{split}
\]
Thus $r_1-r_3=\sqrt2 s(\gamma-\delta)$ and, once
$\gamma=\delta$, $r_2-r_0=4as\alpha$. Column 3 reverses the
signs of these four residuals. This proves the asserted equivalence.

If $a=0$, then $b=\pm1$, and $s\ne0$ forces
$n\cdot p=n\cdot w=0$. Hence $v_0=\pm n$ and $x,y$ are orthogonal
unit vectors in its perpendicular plane, giving class A.
If $a\ne0$, then $n\cdot v_0=0$ and
$n\cdot p=n\cdot w=\varepsilon/\sqrt2$.
The last equation and $b^2+4a^2=1$ give
$a=\eta s/2$ and $b=-\varepsilon\eta c$.
Substitution, using $p=\varepsilon(n+n\times v_0)/\sqrt2$, gives class B.
The exact polynomial verifier independently checks the unit norms and
all triangle equations for all eight signed circles.
\end{proof}

\begin{lemma}[Complete third settings reduce to class A]
\label{lem:commutingcompletion}
For $s\ne0$, every complete third setting is equivalent, allowing an
exchange of the first two settings, to four class-A outcomes with tags
\[
 (1,1,0),\quad(1,1,\pi),\quad(-1,-1,\beta),\quad(-1,-1,\beta+\pi).
\]
The first pair retains the form $U(q)$, possibly replacing $q$ by
$-q$ or $\overline q$.
\end{lemma}
\begin{proof}
For unscaled normalized frames, orthogonality is equivalent to
$\sum_a\overline{u_a}u'_a=0$.
Within either class the squared norm of this sum is zero when both
signs reverse, $8(1+\cos(t'-t))$ when signs coincide, and four when
exactly one sign reverses. Between classes, after rotating the first
angle to zero, it is
\[
 4+L\cos t'+M\sin t',\qquad
 L^2+M^2=8(1+\rho c),\quad \rho\in\{1,-1\}.
\]
Since $s\ne0$, its minimum is strictly positive. Consequently no
complete setting mixes the classes. Within one class, four orthogonal
outcomes consist of two antipodal points on one signed circle and two
on the circle with both signs reversed.

Exchange $P,Q$ and swap their outcome indices 1 and 2. The first-pair
array becomes $U(\overline q)$, rather than $U(q)$; this is the
conjugation of the phase allowed in the lemma statement. For a class-B frame, the new normalized
frame has zero real part in its first nontrivial row and has second
nontrivial row $-\varepsilon\mathbf i$. It therefore belongs to class A.
Explicitly the new frame is $S_b\overline{S_0}$, with row order
$(0,2,1,3)$ and $S_b=\frac12\sum_a\overline{U_{ab}}u_a$.
Equivalently, the new frame has zero real part in its first nontrivial
row, which is the case $a=0$ of Lemma~\ref{lem:commutingoutcomes} for the
array $U(\overline q)$; since $s\neq0$ is preserved, that case gives
class A. The verifier also checks this exchange identity exactly.

In class A choose the signed circle with first sign positive.
If its second sign is negative, interchange computational rows 1 and 3;
this sends $U(q)$ to $U(-q)$ and reverses that sign.
A common internal rotation about $\mathbf i$ sets the first angle to zero.
The second angle $\beta$ remains arbitrary, giving the displayed setting.
\end{proof}

\paragraph{Completion of the proof of Theorem~\ref{thm:commutingnonextension}.}
\begin{proof}
Use Lemma~\ref{lem:commutingpair}. When $s=0$, the first pair is
real-Hadamard and Theorem~\ref{thm:allfixedpair} applies.
Otherwise Lemma~\ref{lem:commutingcompletion} gives the normalized class-A
setting. Any prospective fourth outcome is quaternionic in the same
coordinates, by rigidity applied to its completion of the first pair.
Its deviation $D=T-\id/4$ therefore lies in the 28-dimensional real
quaternionic Hermitian space and satisfies
\[
 \operatorname{tr}\bigl(D\Sym(P_a,Q_b,R_c)\bigr)=0
 \quad\text{for all }a,b,c.
\]

For the Pauli expansion in this proof we use the embedding
\[
 \mathbf i\longmapsto i\sigma_z,\qquad
 \mathbf j\longmapsto i\sigma_y,\qquad
 \mathbf k\longmapsto i\sigma_x.
\]
Appendix~\ref{app:explicit} instead uses
$(\mathbf i,\mathbf j,\mathbf k)\mapsto
(i\sigma_x,-i\sigma_y,i\sigma_z)$.
The two conventions are related by conjugation on the internal qubit
with the Hadamard unitary
$H_2=2^{-1/2}\bigl(\begin{smallmatrix}1&1\\1&-1\end{smallmatrix}\bigr)$:
$H_2\sigma_xH_2=\sigma_z$, $H_2\sigma_zH_2=\sigma_x$, and
$H_2\sigma_yH_2=-\sigma_y$.
Every Pauli word below uses the first convention, with the internal
qubit as its last factor; this simultaneous unitary change preserves
all measurement identities and the obstruction.

Put $d=\cos\beta$, $e=\sin\beta$ and use Walsh observables
for the three settings. For $P$ these are
$A=(III,ZII,IZI,ZZI)$; for $Q$ they are $B_0=III$, $B_1=XII$ and
\[
 B_2=\frac{1+c}{2}IXI+\frac{1-c}{2}XXI
       +\frac{s}{2}(IYZ-XYZ),\quad
 B_3=\frac{1-c}{2}IXI+\frac{1+c}{2}XXI
       -\frac{s}{2}(IYZ-XYZ).
\]
For $R$ they are $C_0=III$, $C_1=YIZ$ and
\[
\begin{gathered}
 C_2=\frac{1+d}{2}ZYY+\frac{1-d}{2}XYX
       +\frac e2(XYY+ZYX),\\
 C_3=\frac{1-d}{2}ZYY+\frac{1+d}{2}XYX
       -\frac e2(XYY+ZYX).
\end{gathered}
\]
These formulas follow by expanding the outcome frames.
The invertible Walsh transform replaces the trace constraints by those
with $\Sym(A_i,B_j,C_k)$.

Grade the quaternionic Hermitian Pauli basis by conjugation with
$J=IZI$. Each of these moment rows has a single grade. The twelve
commuting Pauli coordinates split into eight coordinates
\[
 III,IZI,XII,XZI,YIZ,YZZ,ZII,ZZI
\]
and four coordinates $(YIX,YIY,YZX,YZY)$.
In the Gram matrix of the moment rows' Pauli coefficients, the first
block is $36I_8$, with zero cross terms. The second block is
\[
 G_4=4I_4+8WW^T,\qquad
 W=\begin{pmatrix}
 s&sd&0&0\\0&-se&0&0\\0&-ce&0&e\\-c&-cd&1&d
 \end{pmatrix}.
\]
Exact Pauli multiplication verifies this identity using only
$c^2+s^2=d^2+e^2=1$. Thus the commuting-coordinate Gram matrix is
strictly positive for every parameter. A candidate deviation must have
all twelve commuting coordinates zero, giving $JDJ=-D$.
Consequently $\operatorname{tr}D^3=0$.

For a rank-two projector on $\mathbb C^8$, however,
$D=T-\id/4$ has eigenvalues $3/4$ twice and $-1/4$ six times,
so $\operatorname{tr}D^3=3/4$, a contradiction.
The standard-library verifier
\texttt{verify\_commuting\_\allowbreak pair\_theorem.py} checks the outcome identities,
orthogonality formulas, exchange reduction, observable formulas and
the polynomial Gram factorization.
\end{proof}

\section{Fixed-real-pair classification and nonextension}\label{app:fixedpair}
\subsection{The fixed real pair and its standard completion}\label{app:explicit}
Use zero-based indices $a,b,c\in\{0,1,2,3\}$. Define
\[
H=\begin{pmatrix}1&1&1&1\\1&-1&1&-1\\
1&1&-1&-1\\1&-1&-1&1\end{pmatrix},\qquad
(E_0,E_1,E_2,E_3)=(\id_2,i\sigma_x,-i\sigma_y,i\sigma_z),
\]
where the Pauli matrices are
\[
\sigma_x=\begin{pmatrix}0&1\\1&0\end{pmatrix},\quad
\sigma_y=\begin{pmatrix}0&-i\\i&0\end{pmatrix},\quad
\sigma_z=\begin{pmatrix}1&0\\0&-1\end{pmatrix}.
\]
On $\mathbb C^4\otimes\mathbb C^2$, set
\begin{align*}
P_a&=|a\rangle\langle a|\otimes\id_2,\\
Q_b&=\frac14\sum_{u,v=0}^3 H_{ub}H_{vb}|u\rangle\langle v|\otimes\id_2,\\
R_c&=\frac14\sum_{u,v=0}^3 H_{uc}H_{vc}|u\rangle\langle v|\otimes E_uE_v^*.
\end{align*}
Each measurement consists of four mutually orthogonal rank-two projections
summing to $\id_8$. The triple obeys
\begin{equation}\label{eq:explicittriple}
 A(BC+CB)A=\frac18 A
\end{equation}
for every selected outcome triple and every anchor $A$ among its three
members. These properties follow from the construction in \cite{DF} and
are also reproduced exactly by the accompanying construction checker.


For the displayed triple and
$O\in O(4)$, set $R_c^O=(O\otimes\id_2)R_c(O^T\otimes\id_2)$,
while keeping $P$ and $Q$ fixed. This operation conjugates only the
third setting, so it is not one of the allowed simultaneous equivalences.

\begin{theorem}[Real-orthogonal deformation]\label{thm:rotations}
For every real orthogonal $O$, the triple $(P,Q,R^O)$ is an algebraic
three-fold unbiased triple with four outcomes of rank two on $\mathbb C^8$.

\end{theorem}
\begin{proof}
For a fixed outcome $c$, put
\[
 U_{ac}=\sum_{v=0}^3 O_{av}H_{vc}E_v,
 \qquad V_c=\frac12\begin{pmatrix}U_{0c}\\U_{1c}\\U_{2c}\\U_{3c}\end{pmatrix}.
\]
The two-sided Clifford identities and orthogonality of $O$ imply
\[
 U_{ac}^*U_{a'c}+U_{a'c}^*U_{ac}=2\delta_{aa'}\id_2,
 \qquad
 U_{ac}U_{a'c}^*+U_{a'c}U_{ac}^*=2\delta_{aa'}\id_2.
\]
In particular every $U_{ac}$ is unitary, $V_c^*V_c=\id_2$, and
$R_c^O=V_cV_c^*$. The $R_c^O$ form a PVM because they are a common
unitary conjugate of the original third setting. Their pairwise MUM
relations with $P$ follow from unitarity of each block. If
$q_b=\frac12\sum_a H_{ab}|a\rangle$, then
$V_c^*(q_b\otimes\id_2)=\frac14\sum_aH_{ab}U_{ac}^*$.
The Clifford identities make the product with its adjoint $\id_2/4$,
in either order, proving the MUM relations with $Q$.

For the $P_a$ anchor, the only nonzero diagonal block of
$P_a(Q_bR_c^O+R_c^OQ_b)P_a$ is
\[
 \frac1{16}\sum_v H_{ab}H_{vb}
 (U_{vc}U_{ac}^*+U_{ac}U_{vc}^*)=\frac18\id_2.
\]
Compression to the $Q_b$ frame gives the same expression with the
products in reversed order of the row labels, hence $\id_2/8$.
Finally compression to the $R_c^O$ frame gives
\[
 V_c^*(P_aQ_b+Q_bP_a)V_c
 =\frac1{16}\sum_v H_{ab}H_{vb}
 (U_{ac}^*U_{vc}+U_{vc}^*U_{ac})=\frac18\id_2.
\]
These are all three anchor identities.

\end{proof}

\subsection{Exhaustive normal form with the first pair fixed}

\begin{lemma}[Outcome-frame reduction]\label{lem:outcomeframe}
Keep the computational $P$ and the real-Hadamard $Q$ of Appendix~\ref{app:explicit} fixed.
Every rank-two outcome $S$ of a third setting in an algebraic three-fold
unbiased triple has the form
\[
 S=(\id_4\otimes V)R_0^{\epsilon}(\id_4\otimes V^*),
 \qquad V\in SU(2),\quad \epsilon\in\{+,-\},
\]
where $R_0^+=R_0$ and $R_0^-$ is obtained by replacing $E_3$ by $-E_3$.
\end{lemma}
\begin{proof}
Write $S_{av}$ for its $2\times2$ row blocks. Pairwise unbiasedness
gives $S_{aa}=\id_2/4$. The $P_a$-anchored triple equation gives, for
each $b$,
\[
 \frac14\sum_v H_{ab}H_{vb}(S_{va}+S_{av})=\frac18\id_2.
\]
The $v=a$ term is already $\id_2/8$. Inverting the Hadamard transform
in $b$ shows $S_{va}+S_{av}=0$ for $v\ne a$. Thus each off-diagonal
block is anti-Hermitian; this conclusion does not impose a quaternionic
ansatz on $S$.

Because $S$ is a rank-two projection, choose an isometry $W$ with
$S=WW^*$. Write its row blocks as $W_a=U_a/2$; the scalar diagonal
blocks show that every $U_a$ is unitary. Off-diagonal anti-Hermiticity
gives $U_aU_v^*+U_vU_a^*=0$. A common right multiplication by $U_0^*$
leaves $S$ unchanged and sets $U_0=\id_2$. Now $U_j^*=-U_j$ and
$U_jU_k+U_kU_j=0$ for distinct $j,k\in\{1,2,3\}$.
None of the $U_j$ can be $\pm i\id_2$, since that would contradict
anticommutation with the other unitary blocks. The Hermitian involutions
$-iU_j$ therefore have eigenvalues $+1,-1$, and can be written
$\boldsymbol n_j\cdot\boldsymbol\sigma$ with real unit vectors
$\boldsymbol n_j$. Anticommutation says the three vectors are orthogonal.
There are two orientations of this ordered orthonormal triad. The
adjoint action of $SU(2)$ realizes every proper rotation, hence the triad
is conjugate to either the explicit $E_1,E_2,E_3$ or
$E_1,E_2,-E_3$. This proves the claimed form.
\end{proof}

\begin{theorem}[Fixed-pair classification]\label{thm:fixedpair}
For the computational/real-Hadamard pair $(P,Q)$ of Appendix~\ref{app:explicit}, a
four-outcome rank-two PVM $S$ completes an algebraic three-fold unbiased
triple if and only if, after relabelling its outcomes,
\[
 S_c=(O\otimes V)R_c(O^T\otimes V^*),
 \qquad O\in O(4),\quad V\in SU(2).
\]
The parameters need not be unique.
\end{theorem}
\begin{proof}
The forward construction follows from Theorem~\ref{thm:rotations}:
$\id_4\otimes V$ commutes with all $P_a,Q_b$, so its conjugation
preserves the triple equations. We prove exhaustiveness.

Introduce two triples of operators on $\mathbb C^4$, with $XY=\sigma_x\otimes\sigma_y$ and similarly for the other
two-letter products, where $I=\id_2$:
\[
 (L_1^+,L_2^+,L_3^+)=(IY,YZ,YX),\qquad
 (L_1^-,L_2^-,L_3^-)=(ZY,YI,-XY).
\]
Both obey the usual Pauli multiplication relations. Their generators
are purely imaginary skew matrices; the two triples commute with one
another and satisfy $\operatorname{tr}(L_j^+L_k^-)=0$. Direct expansion
gives
\[
 R_0^{\pm}=\frac14\bigl(\id_8+
 L_1^{\pm}\otimes\sigma_x-L_2^{\pm}\otimes\sigma_y+
 L_3^{\pm}\otimes\sigma_z\bigr).
\]
Consequently, for arbitrary $V,W\in SU(2)$,
\[
 \operatorname{tr}\left[
 (\id_4\otimes V)R_0^+(\id_4\otimes V^*)
 (\id_4\otimes W)R_0^-(\id_4\otimes W^*)\right]=\frac12.
\]
All cross terms vanish by the stated traces; the identity contribution
is $8/16$. Outcomes of different signs therefore cannot be orthogonal.
Lemma~\ref{lem:outcomeframe} implies that all four outcomes use the
same sign. These finite Pauli identities and the 81 cross-component
trace identities are independently checked in
\texttt{verify\_fixedpair\_structure.py}.

For either sign, the algebra generated by the $L_j^{\pm}$ is a copy
of $M_2(\mathbb C)$ acting with multiplicity two on $\mathbb C^4$.
This follows directly by diagonalizing $L_3^{\pm}$ and using the
anticommuting involution $L_1^{\pm}$ to identify its two eigenspaces.
Thus in suitable factor coordinates each outcome is the identity on
the multiplicity factor tensored with a rank-one maximally entangled
two-qubit projector. The four orthogonal outcomes become an orthonormal
basis of maximally entangled vectors on the two active qubits.

For completeness, the local-unitary normal form of such a basis is
elementary. Write the vectors as
$(\id_2\otimes B_c)|\Phi\rangle$ with unitary $B_c$ and
$|\Phi\rangle=(|00\rangle+|11\rangle)/\sqrt2$.
Orthonormality means $\operatorname{tr}(B_c^*B_d)=2\delta_{cd}$.
A common left multiplication sets $B_0=\id_2$. For $j=1,2,3$, a
traceless unitary $B_j$ is, after changing an irrelevant phase,
$\boldsymbol n_j\cdot\boldsymbol\sigma$ with a real unit vector.
The three vectors are orthonormal. Changing one vector's sign if needed
and conjugating by a common $C\in SU(2)$ reduces them to the Pauli
axes. On the entangled vectors this conjugation is realized by
$\overline C\otimes C$, which preserves $|\Phi\rangle$.
This proves that every such basis is locally equivalent to the Bell basis.
This two-qubit fact is standard; the broader classification of
iso-entangled bases is studied in \cite{DS}.

An arbitrary unitary on the active first qubit is generated, up to an
irrelevant scalar phase, by the $L_j^{\pm}$. Its representatives
$\exp(i\theta\boldsymbol n\cdot\boldsymbol L^{\pm})$ are real
orthogonal $4\times4$ matrices, because the $L_j^{\pm}$ are purely
imaginary skew. The local unitary on the second active qubit is the
factor $V$. The original four projectors $R_c^{\pm}$ already form
the Bell basis in these factor coordinates: their Pauli correlation
signs are the four triples with product $-1$. We therefore obtain
the displayed normal form with the common sign fixed. Finally,
$D=\operatorname{diag}(1,1,1,-1)$ satisfies
$R_c^-=(D\otimes\id_2)R_c^+(D\otimes\id_2)$, so a negative sign
can be absorbed into $O$.
\end{proof}

The deformation theorem supplies completions of the fixed real pair.
Theorem~\ref{thm:fixedpair} proves that these exhaust its complete third
settings, and the next theorem excludes their four-fold extensions.
These fixed-pair statements alone do not classify triples with an
arbitrary first pair.
Its proof uses standard Clifford and pair-normal-form ingredients.

\subsection{Nonextension for every completion of the fixed pair}
\begin{theorem}[All fixed-pair completions are nonextendible]
\label{thm:allfixedpair}
Every four-outcome rank-two third setting $S$ such that $(P,Q,S)$ is a
3-UM, with the computational/real-Hadamard pair of Appendix~\ref{app:explicit} fixed, has
no algebraic four-fold unbiased extension. This includes every member
of the orthogonal family in Theorem~\ref{thm:fixedpair}, without a genericity
condition or an exclusion of exceptional parameters.
\end{theorem}
\begin{proof}
Use $L^+=(IY,YZ,YX)$ and $L^-=(ZY,YI,-XY)$ from the fixed-pair proof.
An outcome of any third setting is a maximally entangled projector on
one of these spin factors and the internal qubit. Consequently its
centered part has the form
\[
 T-\id_8/4=\frac14\sum_{k,l=1}^3 X_{kl}L_k^\epsilon\otimes\sigma_l,
 \qquad X\in O(3),\quad \det X=-1.
\]
This is the usual orthogonal correlation matrix of a maximally entangled
two-qubit pure state. In particular its three columns are nonzero.

Normalize the given third setting to the positive type. Swapping the two
first-factor qubits exchanges $L^+$ and $L^-$ up to signed permutations;
it permutes the computational and real-Hadamard outcomes, so preserves
the fixed pair as measurements. The common internal unitary can also be
removed. In these coordinates the third outcomes have the form
\[
 S_c=\frac14\left(\id_8+
 \sum_{i,j=1}^3 M_{ij}s_{jc}L_i^+\otimes\sigma_j\right),
 \qquad M\in SO(3),
\]
where the four sign vectors $s_{\cdot c}$ are the Bell correlation signs.
They satisfy $\sum_c s_{jc}=0$ and
$\sum_c s_{jc}s_{kc}=4\delta_{jk}$. Thus taking the corresponding
signed sums of the $P$-anchored four-fold equations isolates the operator
$\sum_i M_{ij}L_i^+\otimes\sigma_j$ for each $j$.

Write $m=(M_{1j},M_{2j},M_{3j})$ and
$x=(X_{1j},X_{2j},X_{3j})$. Exact Pauli expansion of these equations
gives, in the same-type case,
\begin{equation}\label{eq:columnobstruction}
 m\cdot x=0,\qquad m_i x_k+m_k x_i=0\quad(i\ne k).
\end{equation}
For clarity, one may extract these four equations using only anchor
$P_0$ and the four $Q_b$: up to a common nonzero scale their scalar
coefficients are $3m\cdot x$ plus the three off-diagonal expressions,
with coefficient sign rows
$(1,1,1),(-1,1,-1),(-1,-1,1),(1,-1,-1)$.
These four rows form an invertible Hadamard transform together with the
constant column. In the opposite-type case, the equations with all four
$P_a$ anchors give instead
\begin{equation}\label{eq:oppositecolumn}
 m_i x_k=0\qquad\hbox{for all }i,k.
\end{equation}
Both coefficient identities hold symbolically for all real $M,X$;
\texttt{verify\_fixedpair\_no\_extension.py} independently checks their
row spans by exact Gaussian-integer Pauli multiplication and rational
row reduction. It checks all 27 monomials per column, for both types,
so no parameter sampling enters this reduction.

Since $m$ is a unit vector, \eqref{eq:oppositecolumn} forces $x=0$,
impossible for an orthogonal $X$. Consider \eqref{eq:columnobstruction}.
If $m_i\ne0$ and $m_k=0$, then $x_k=0$. On the support of $m$, put
$r_i=x_i/m_i$. The off-diagonal equations say $r_i+r_k=0$.
For three supported coordinates this forces all $r_i=0$; for one
coordinate the dot equation forces $x=0$. A nonzero $x$ is therefore
possible only when $m$ has exactly two supported coordinates $i,k$,
with $|m_i|=|m_k|=1/\sqrt2$ and $r_i=-r_k$.

An orthogonal $3\times3$ matrix cannot have three columns of that form.
Columns with different two-coordinate supports overlap in exactly one
coordinate and have dot product $\pm1/2$; columns with identical support
lie in a two-dimensional plane and cannot supply three orthogonal vectors.
At least one column of $M$ therefore forces the corresponding column of
$X$ to vanish, again a contradiction. Every prospective fourth outcome
must also complete a 3-UM with $P,Q$ by heredity, so the argument excludes
every possible fourth setting.
\end{proof}

The theorem is a fixed-pair result, not a claim that an arbitrary pair
in a strong triple is equivalent to the computational/real-Hadamard pair.

\section{An auxiliary complete-frame formulation}\label{app:noncommuting}
\subsection{Strong settings from supported kernels}
\label{sec:globalcoarse}
This reduction includes arbitrary outcome classes and every parameter
stratum. Write $\tau(A)=\operatorname{Re}\operatorname{Tr}_{\mathbb H}A$.
For a fixed quaternionic MUM pair $P,Q$, define the real linear space
\[
 K_{P,Q}=\{D\in\operatorname{Herm}_4(\mathbb H):
 \tau((P_aQ_b+Q_bP_a)D)=0\text{ for every }a,b\}.
\]
Its orthogonal complement contains $I$, every $P_a$ and every $Q_b$.
Thus elements of $K_{P,Q}$ are traceless and have zero diagonal in
both measurement bases. No dimension or constant-rank assumption is made.

\begin{proposition}[Complete settings from supported kernels]
\label{prop:globalkernels}
A quaternionic rank-one projector $T$ completes $P,Q$ to a strong
outcome triple precisely when $T-I/4\in K_{P,Q}$.
Complete third settings are equivalently obtained as follows. Choose
$X\in K_{P,Q}$ with $X^2=I$, set $e_\pm=(I\pm X)/2$, and define
\[
 L_\pm(X)=\{A\in K_{P,Q}:XA=\pm A\}.
\]
A complete setting with this balanced coarse-graining exists if and only
if both $L_+(X)$ and $L_-(X)$ are nonzero.
\end{proposition}
\begin{proof}
Quaternionic rank-one compression of a Hermitian matrix is scalar.
For a MUM pair, $\tau(P_aQ_b)=1/4$. The strong identities are therefore
exactly $\tau((P_aQ_b+Q_bP_a)T)=1/8$; summing these equations also
gives both pairwise unbiasedness relations. Subtracting $I/4$ gives
membership in $K_{P,Q}$. Cyclicity of the real trace gives all three
anchors, including the anchor at $T$.

Given a complete setting $R_0,\ldots,R_3$, put
$X=R_0+R_1-R_2-R_3$. Then $X\in K_{P,Q}$ and $X^2=I$,
while $R_0-R_1\in L_+(X)$ and $R_2-R_3\in L_-(X)$ are nonzero.
Conversely, $X$ is traceless, so each $e_\pm$ has quaternionic rank two.
Hermiticity gives $AX=XA=\pm A$ for $A\in L_\pm(X)$.
Such $A$ is supported on $e_\pm$ and traceless. A traceless Hermitian
$2\times2$ quaternionic matrix squares to a nonnegative scalar multiple
of the identity: in coordinates
$\bigl(\begin{smallmatrix}a&q\\\bar q&-a\end{smallmatrix}\bigr)^2
=(a^2+|q|^2)I$. Normalize nonzero choices to
$A_\pm^2=e_\pm$. Then
\[
 (e_++A_+)/2,\quad(e_+-A_+)/2,\quad
 (e_-+A_-)/2,\quad(e_--A_-)/2
\]
are four orthogonal quaternionic rank-one projectors summing to $I$.
Each differs from $I/4$ by an element of $K_{P,Q}$, hence the resulting
complete setting is strong. This also proves exhaustion.
\end{proof}

\section{Verification artifacts}\label{app:verification}
The source and every verifier below, together with the shared module
\texttt{helpers.py} and the helper scripts it relies on, are available at
\url{https://github.com/danielrupawalla/rank-two-obstruction}.
The repository includes a README, a dependency declaration and a runner
that executes each verifier with assertions enabled. The core moment and commuting-pair verifiers use Python's standard
library. The polynomial chart and global-classification checks use SymPy.
\small
\begin{longtable}{p{0.43\textwidth}p{0.51\textwidth}}
\toprule Result & Verifier or artifact\\\midrule
Stratum symmetries and embedding change & \texttt{verify\_stratum\_symmetry.py}\\
Global outcome and setting identities & \texttt{verify\_global\_outcome\_classification.py}\\
Global linear and exceptional reductions & \texttt{verify\_global\_linear\_reduction.py}, \texttt{verify\_global\_exceptional\_reduction.py}, \texttt{verify\_global\_exceptional\_complex\_term.py}\\
Global complete-frame reduction & \texttt{verify\_global\_coarse\_reduction.py}\\
Unified noncommuting pair chart & \texttt{verify\_unified\_pair\_chart.py}\\
Coarse-support bridge and two-column chart & \texttt{verify\_coarse\_support\_bridge.py}, \texttt{verify\_two\_column\_kernel.py}\\
Spectral moment rigidity & \texttt{verify\_spectral\_moment\_rigidity.py}\\
Quaternionic triangle normal form & \texttt{verify\_triangle\_normal\_form.py}\\
Scalar-triangle exclusion & \texttt{verify\_scalar\_triangle\_exclusion.py}\\
Uniform commuting-pair theorem & \texttt{verify\_commuting\_pair\_theorem.py}\\
Fixed-real-pair structure and obstruction & \texttt{verify\_fixedpair\_structure.py}, \texttt{verify\_fixedpair\_no\_extension.py}\\
\bottomrule
\end{longtable}
\normalsize
The spectral verifier checks symbolic identities in $k$ and $1/n$, all
relevant fourth- and sixth-order word types, and exact matrix examples
with positive defects. The global-classification verifiers check the
symbolic matrix identities and branch reductions used in the analytic proof.

\end{document}
